\documentclass{uai2026} % anonymous submission

\usepackage[american]{babel}
\usepackage{natbib}
    \bibliographystyle{plainnat}
    \renewcommand{\bibsection}{\subsubsection*{References}}
\usepackage{mathtools}
\usepackage{booktabs}
\usepackage{tikz}
\usepackage{algorithm2e}
\usepackage{amssymb}
\usepackage{multirow}
\usepackage{xcolor}
\usepackage{amsthm}

% ---- Theorems ----
\newtheorem{theorem}{Theorem}[section]
\newtheorem{proposition}[theorem]{Proposition}
\newtheorem{lemma}[theorem]{Lemma}
\newtheorem{corollary}[theorem]{Corollary}
\theoremstyle{definition}
\newtheorem{definition}[theorem]{Definition}
\theoremstyle{remark}
\newtheorem{remark}[theorem]{Remark}

% ---- Macros ----
\newcommand{\RN}{\mathrm{RN}}    % Ratio Norm
\newcommand{\bz}{\mathbf{z}}
\newcommand{\bg}{\mathbf{g}}
\newcommand{\bu}{\mathbf{u}}
\newcommand{\bv}{\mathbf{v}}
\newcommand{\bw}{\mathbf{w}}
\newcommand{\bx}{\mathbf{x}}
\newcommand{\by}{\mathbf{y}}
\newcommand{\bq}{\mathbf{q}}
\newcommand{\bd}{\mathbf{d}}
\newcommand{\be}{\mathbf{e}}
\newcommand{\bW}{\mathbf{W}}
\DeclareMathOperator{\prox}{prox}
\DeclareMathOperator{\soft}{soft}
\DeclareMathOperator{\sgn}{sgn}
\DeclareMathOperator*{\argmin}{arg\,min}

\title{Linearized ADMM with Ratio Norm for Neural Network Feature Selection}

\author[1]{Anonymous Author(s)}
\affil[1]{Anonymous Institution}

\begin{document}
\maketitle

% ================================================================
%  ABSTRACT
% ================================================================
\begin{abstract}
Feature selection in high-dimensional data with complex nonlinear dependencies remains challenging.
Existing neural network approaches either apply soft gating with standard sparsity penalties, yielding diffuse feature scores, or require specialized architectures that limit flexibility.
We propose a simple yet effective framework combining a \emph{Gated Feature-Selection MLP} with \emph{Linearized ADMM} and a \emph{Ratio Norm} ($\lVert\bz\rVert_1/\lVert\bz\rVert_2$) penalty.
The gated network multiplies each input feature by a learnable scalar gate; during training, the gate vector is optimized jointly with network weights via Adam, while the ADMM auxiliary variable receives an epoch-level proximal update under the scale-invariant Ratio Norm.
Feature dropout at 70\% prevents memorization through noise features.
On synthetic benchmarks with up to 2{,}048 dimensions, our method achieves 66.8\% average best-$k$ accuracy, ranking first among all neural network methods and third overall.
On a causal DAG dataset (2{,}000 features), it attains the highest fork-feature recovery (32.4\% best-$2k_2$), substantially outperforming TreeSHAP (17.3\%).
On the MADELON real-world benchmark, it matches Random Forest and TreeSHAP (AUROC 0.964), exceeding all other neural approaches.
These results demonstrate that Linearized ADMM with Ratio Norm provides a principled and competitive approach to neural feature selection, especially for detecting nonlinear and interaction-based informative features.
\end{abstract}

% ================================================================
%  1. INTRODUCTION
% ================================================================
\section{Introduction}\label{sec:intro}

Feature selection---identifying a small subset of informative variables from a potentially large candidate set---is a fundamental problem in machine learning \citep{guyon2003introduction, li2018feature}.
In modern high-dimensional settings, the challenge is compounded by nonlinear feature interactions: purely univariate statistics (e.g., mutual information) fail to detect XOR-like patterns where individual features carry no marginal information, yet jointly determine the response.

Traditional filter methods such as mutual information \citep{peng2005mrmr} and Relief \citep{kira1992relief} scale well but are limited to detecting marginal or pairwise associations.
Tree-based methods \citep{breiman2001random, lundberg2017shap} capture nonlinear relationships through ensemble averaging, but their importance scores reflect split-based relevance rather than a principled sparsity mechanism.
Neural network approaches---LassoNet \citep{lemhadri2021lassonet}, CancelOut \citep{borisov2019cancelout}, Concrete Autoencoders \citep{balin2019concrete}, and DeepPINK \citep{lu2018deeppink}---embed feature selection within the learning pipeline, but often rely on $\ell_1$ penalties or Gumbel-Softmax relaxations that are sensitive to feature scale and provide limited theoretical guarantees on exact sparsity.

We propose a framework consisting of three components:
\begin{enumerate}
    \item \textbf{Gated Feature-Selection MLP}: A lightweight two-layer network with an unbounded learnable gate vector $\bg \in \mathbb{R}^m$ applied as element-wise multiplication, combined with 70\% feature dropout during training.
    \item \textbf{Linearized ADMM}: The gate vector participates in standard Adam optimization (preserving full gradient dynamics), while an auxiliary variable $\bz$ is updated at the epoch level through a closed-form proximal step.
    \item \textbf{Ratio Norm penalty}: The scale-invariant sparsity measure $R(\bz) = \lVert\bz\rVert_1 / \lVert\bz\rVert_2$ is used in the $\bz$-step, ensuring that features with different magnitudes are treated fairly.
\end{enumerate}

Our contributions are:
\begin{itemize}
    \item A principled integration of ADMM-based structured sparsity with neural network training that maintains the gate within the optimizer's gradient flow.
    \item A proximal operator for the Ratio Norm penalty, computed via a cubic equation with closed-form solution.
    \item Comprehensive experiments showing state-of-the-art performance among neural methods on synthetic (nonlinear/interaction), causal (DAG), and real-world (MADELON) benchmarks.
\end{itemize}

% ================================================================
%  2. RELATED WORK
% ================================================================
\section{Related Work}\label{sec:related}

\paragraph{Filter and tree-based methods.}
Univariate mutual information (MI) ranks features independently, while mRMR \citep{peng2005mrmr} accounts for inter-feature redundancy.
Relief \citep{kira1992relief} uses nearest-neighbor distances.
Random Forests \citep{breiman2001random} provide Gini importance; TreeSHAP \citep{lundberg2017shap} yields theoretically grounded Shapley-value importances but inherits the tree ensemble's limitations in extremely high dimensions.

\paragraph{Neural feature selection.}
LassoNet \citep{lemhadri2021lassonet} introduces hierarchical $\ell_1$ penalties on skip connections.
CancelOut \citep{borisov2019cancelout} applies sigmoid- or softmax-gated masks.
Concrete Autoencoders \citep{balin2019concrete} use Gumbel-Softmax for differentiable discrete selection.
Stochastic Gates \citep{yamada2020stg} employ Gaussian-smoothed Bernoulli masks.
DeepPINK \citep{lu2018deeppink} combines knockoff statistics with pairwise feature layers.
All these methods couple the sparsity mechanism tightly with the gradient update, making it difficult to enforce hard sparsity or use non-smooth penalties beyond $\ell_1$.

\paragraph{ADMM in deep learning.}
ADMM has been applied to neural network weight pruning \citep{ren2019admm} and training without gradients \citep{taylor2016admm}.
Linearized ADMM variants \citep{lin2011linearized, yang2013linearized} have been studied for convex problems.
\citet{yu2026sadmm} proposed Scope-Driven SADMM, which bridges ADMM with stochastic gradient descent for network pruning via a three-variable splitting of the Ratio Norm.
Our work adapts this framework specifically for the \emph{feature gate vector} in a neural network for feature selection, with importance-adaptive regularization.

\paragraph{Ratio Norm and sparse penalties.}
The $\ell_1/\ell_2$ ratio was studied as a sparseness measure by \citet{hoyer2004sparseness} in the context of non-negative matrix factorization.
\citet{nasery2024end} applied the ratio norm to end-to-end sparse model training.
Unlike the standard Lasso \citep{tibshirani1996lasso}, which penalizes each coefficient independently and is scale-sensitive, the Ratio Norm is invariant to global scaling of $\bz$, making it especially suitable when gate magnitudes vary across features.

% ================================================================
%  3. METHOD
% ================================================================
\section{Method}\label{sec:method}

\subsection{Gated Feature-Selection MLP}\label{sec:model}

Given input $\bx \in \mathbb{R}^m$, we define a learnable gate vector $\bg \in \mathbb{R}^m$ initialized to $\mathbf{1}$, and compute:
\begin{equation}\label{eq:gating}
    \hat{\bx} = \bx \odot (\bg \odot \mathbf{d}),
\end{equation}
where $\mathbf{d} \sim \text{Bernoulli}(1 - p)^m / (1-p)$ is an inverted-dropout mask with drop rate $p = 0.7$ applied during training.
The masked input passes through a two-layer MLP:
\begin{equation}\label{eq:mlp}
    h = \phi(\bW_2\, \phi(\bW_1\, \hat{\bx} + b_1) + b_2), \quad y = \bW_3\, h + b_3,
\end{equation}
where $\phi$ is the Mish activation \citep{misra2019mish} and hidden dimensions are 32.

\paragraph{Design choices.}
The gate is \emph{unbounded}---unlike CancelOut's sigmoid gating $\sigma(\theta) \in [0,1]$---allowing the ADMM penalty to drive entries to exact zero.
The high feature dropout rate of 70\% is critical: it prevents the network from memorizing through noise features, forcing the gate to concentrate weight on truly informative features.
The compact architecture (2$\times$32 vs.\ 5$\times$58 in the baseline MLP) complements the dropout by reducing model capacity.

\subsection{Three-Variable ADMM Splitting for the Ratio Norm}\label{sec:admm}

The Ratio Norm $R(\bz) = \lVert\bz\rVert_1 / \lVert\bz\rVert_2$ couples all components through both the numerator and denominator, preventing a direct proximal operator.
Following the SADMM framework \citep{yu2026sadmm}, we introduce \emph{two} auxiliary copies of the gate vector---$\by$ to handle the $\ell_2$ denominator and $\bz$ to handle the $\ell_1$ numerator---yielding the constrained problem:
\begin{equation}\label{eq:admm-obj}
    \min_{\theta,\bg}\;\mathcal{L}(\theta,\bg) + \frac{C}{N}\frac{\lVert\bz\rVert_1}{\lVert\by\rVert_2}
    \quad\text{s.t.}\;\bg = \by,\;\bg = \bz,
\end{equation}
where $\theta = \{\bW_i, b_i\}$ and $\lambda_j = C / s_j$ with $s_j = \lVert\bW_1[:,j]\rVert_2 + \epsilon$.
The augmented Lagrangian in scaled form is:
\begin{equation}\label{eq:auglag}
    \mathcal{L}_\rho = \mathcal{L}(\theta,\bg)
    + \frac{C}{N}\frac{\lVert\bz\rVert_1}{\lVert\by\rVert_2}
    + \frac{\rho}{2}\lVert\bg - \by + \bv/\rho\rVert^2
    + \frac{\rho}{2}\lVert\bg - \bz + \bw/\rho\rVert^2,
\end{equation}
where $\bv,\bw$ are dual variables.

\paragraph{Training proceeds in two phases:}

\textbf{Phase 1 (Warm-up, epochs $0$ to $T_w{-}1$).}
Standard Adam \citep{kingma2015adam} training on $\mathcal{L}(\theta,\bg)$ with feature dropout, no sparsity.

\textbf{Phase 2 (ADMM, epochs $T_w$ to $T{-}1$).}
At each epoch, four subproblems are solved:

\textit{Step 1 ($\bg$-step):} Minimizing~\eqref{eq:auglag} over $(\theta,\bg)$ while fixing $\by,\bz,\bv,\bw$.
Taking the derivative and setting to zero:
\begin{equation}\label{eq:gstep}
    \bg^{(k+1)} = \tfrac{1}{2}\!\left(\by^{(k)} + \bz^{(k)}
    - \tfrac{\bv^{(k)}}{\rho} - \tfrac{\bw^{(k)}}{\rho}
    - \tfrac{1}{\rho}\nabla_{\bg}\mathcal{L}\big|_{\bg^{(k+1)}}\right).
\end{equation}
Since computing the gradient at $\bg^{(k+1)}$ is implicit, we approximate it using Adam's momentum estimate $\hat{m}_k$ (SADMM, see Section~\ref{sec:sadmm}).
Crucially, $\bg$ remains within Adam's state, preserving full gradient dynamics.

\textit{Step 2 ($\by$-step):} With $c^{(k)} = \lVert\bz^{(k)}\rVert_1$ and $\bd^{(k)} = \bg^{(k+1)} + \bv^{(k)}/\rho$, we solve (see Section~\ref{sec:ystep}):
\begin{equation}\label{eq:ysubproblem}
    \by^{(k+1)} = \argmin_{\by}\;\frac{C\, c^{(k)}}{N\lVert\by\rVert_2}
    + \frac{\rho}{2}\lVert\by - \bd^{(k)}\rVert^2.
\end{equation}

\textit{Step 3 ($\bz$-step):} Soft-thresholding (see Section~\ref{sec:zstep}):
\begin{equation}\label{eq:zstep}
    \bz^{(k+1)} = S\!\left(\bg^{(k+1)} + \tfrac{\bw^{(k)}}{\rho},\;
    \tfrac{C}{N\rho\lVert\by^{(k+1)}\rVert_2}\right).
\end{equation}

\textit{Step 4 (Dual updates):}
$\bv^{(k+1)} = \bv^{(k)} + \rho(\bg^{(k+1)} - \by^{(k+1)})$,
$\bw^{(k+1)} = \bw^{(k)} + \rho(\bg^{(k+1)} - \bz^{(k+1)})$.

Additionally, $s_j$ is refreshed each epoch, and $\rho$ is adapted every 5 epochs via residual-balancing \citep{boyd2011admm} with dual rescaling $\bv \leftarrow \bv \cdot \rho_\text{old}/\rho_\text{new}$, $\bw \leftarrow \bw \cdot \rho_\text{old}/\rho_\text{new}$.

\subsection{Solving the $\by$-Subproblem: Cubic Equation}\label{sec:ystep}

The $\by$-subproblem~\eqref{eq:ysubproblem} handles the $\ell_2$ denominator of the Ratio Norm.
Taking the derivative and setting to zero:
\begin{equation}\label{eq:yderiv}
    -\frac{C\, c^{(k)}}{N\lVert\by\rVert_2^3}\by + \rho(\by - \bd^{(k)}) = \mathbf{0}.
\end{equation}
When $\bd^{(k)} \neq \mathbf{0}$, we write $\by^{(k+1)} = \tau^{(k)}\bd^{(k)}$ for a scalar $\tau > 0$.
Substituting into~\eqref{eq:yderiv} and simplifying:
\begin{equation}\label{eq:cubic}
    (\tau^{(k)})^3 - (\tau^{(k)})^2 - D^{(k)} = 0,
\end{equation}
where
\begin{equation}\label{eq:D}
    D^{(k)} = \frac{C\,c^{(k)}}{N\rho\lVert\bd^{(k)}\rVert_2^3}.
\end{equation}
This cubic has a unique positive real root given by Cardano's formula:
\begin{equation}\label{eq:tau}
    \tau^{(k)} = \tfrac{1}{3} + \tfrac{1}{3}\bigl(\mathcal{C}^{(k)} + 1/\mathcal{C}^{(k)}\bigr),
\end{equation}
where
\begin{equation}\label{eq:cardano}
    \mathcal{C}^{(k)} = \sqrt[3]{\frac{27D^{(k)} + 2 + \sqrt{(27D^{(k)}+2)^2 - 4}}{2}}.
\end{equation}
When $\bd^{(k)} = \mathbf{0}$, we set $\by^{(k+1)} = \be^{(k)}$, a random vector with $\lVert\be^{(k)}\rVert_2 = (c^{(k)}/\rho)^{1/3}$.

\subsection{Solving the $\bz$-Subproblem: Soft-Thresholding}\label{sec:zstep}

The $\bz$-subproblem handles the $\ell_1$ numerator:
\begin{equation}\label{eq:zsub}
    \bz^{(k+1)} = \argmin_{\bz}\;\frac{C}{N\lVert\by^{(k+1)}\rVert_2}\lVert\bz\rVert_1
    + \frac{\rho}{2}\Big\lVert\bg^{(k+1)} - \bz + \tfrac{\bw^{(k)}}{\rho}\Big\rVert^2.
\end{equation}
Differentiating and applying the standard soft-thresholding operator $S(\cdot,\cdot)$:
\begin{equation}\label{eq:zupdate}
    z_j^{(k+1)} = S\!\left(g_j^{(k+1)} + \frac{w_j^{(k)}}{\rho},\;\frac{C}{N\rho\lVert\by^{(k+1)}\rVert_2}\right),
\end{equation}
where $S(a,\kappa) = \sgn(a)\max(|a|-\kappa, 0)$.

\begin{proposition}[Pruning Bound]\label{prop:pruning}
Feature $j$ is pruned (i.e., $z_j^{(k+1)} = 0$) if and only if
\begin{equation}\label{eq:prune-bound}
    \left|g_j^{(k+1)} + \frac{w_j^{(k)}}{\rho}\right| \leq \frac{C}{N\rho\lVert\by^{(k+1)}\rVert_2}.
\end{equation}
The pruning threshold adapts automatically: as the $\ell_2$ norm $\lVert\by\rVert_2$ grows (more features active), the threshold shrinks, making pruning more selective.
\end{proposition}

\subsection{Connection to Stochastic Optimization (SADMM)}\label{sec:sadmm}

The $\bg$-step~\eqref{eq:gstep} requires the implicit gradient $\nabla_{\bg}\mathcal{L}|_{\bg^{(k+1)}}$.
Following \citet{yu2026sadmm}, we bridge ADMM with stochastic gradient descent by setting:
\begin{equation}\label{eq:sadmm-bridge}
    \frac{1}{\rho}\nabla_{\bg}\mathcal{L}\big|_{\bg^{(k+1)}} \approx \alpha_k \hat{m}_k,
    \quad \rho = \frac{1}{\alpha_k},
\end{equation}
where $\alpha_k = \eta / (\sqrt{\hat{v}_k} + \epsilon)$ is Adam's adaptive step size and $\hat{m}_k$ is the bias-corrected first moment.
This yields the practical update:
\begin{equation}\label{eq:g-practical}
    \bg^{(k+1)} = \tfrac{1}{2}\bigl(\by^{(k)} + \bz^{(k)}
    - \alpha_k\bv^{(k)} - \alpha_k\bw^{(k)} - \alpha_k\hat{m}_k\bigr).
\end{equation}
This formulation makes the framework compatible with any gradient-based optimizer (SGD, Adam, AdaGrad) and enables mini-batch training without requiring closed-form solutions to the $\bg$-subproblem.

\begin{proposition}[$\ell_1$ Degeneration]\label{prop:l1}
When the Ratio Norm $\lVert\bz\rVert_1/\lVert\bz\rVert_2$ is replaced by the $\ell_1$ norm $\lVert\bz\rVert_1$, the three-variable split reduces to a two-variable split ($\bg = \bz$), and the $\bz$-step becomes:
\begin{equation}
    z_j^{(k+1)} = S\bigl(g_j^{(k)} - \alpha_k\nabla_{g_j}\mathcal{L},\; C\alpha_k/N\bigr),
\end{equation}
which recovers the classical proximal gradient (iterative soft-thresholding) method \citep{parikh2014proximal}.
\end{proposition}

\subsection{Ratio Norm Properties}\label{sec:rationorm}

\begin{definition}[Ratio Norm]
For $\bz \neq \mathbf{0}$, the Ratio Norm is $R(\bz) = \lVert\bz\rVert_1 / \lVert\bz\rVert_2$.
It is scale-invariant: $R(\alpha\bz) = R(\bz)$ for $\alpha > 0$.
Its range is $[1, \sqrt{m}]$, with $R(\bz)=1$ iff at most one component is nonzero (maximally sparse) and $R(\bz) = \sqrt{m}$ iff all components are equal \citep{hoyer2004sparseness}.
\end{definition}

Compared to $\ell_1$, the Ratio Norm has two key advantages for feature selection:
(1)~\emph{Scale invariance}: features with different gate magnitudes are treated equally, avoiding the bias where $\ell_1$ over-penalizes large-magnitude gates;
(2)~\emph{Decoupled sparsity and scale}: the three-variable split~\eqref{eq:admm-obj} separates the sparsity-inducing $\ell_1$ numerator (handled by $\bz$) from the normalizing $\ell_2$ denominator (handled by $\by$), enabling independent closed-form solutions.

The importance-adaptive penalty $\lambda_j = C/s_j$ assigns weaker regularization to features connected to high-norm first-layer weights, automatically calibrating the sparsity threshold.

\subsection{Algorithm Summary}

The complete algorithm is given as Algorithm~\ref{alg:main}.

\begin{algorithm}[t]
\caption{Linearized ADMM with Ratio Norm (SADMM-FS)}\label{alg:main}
\SetAlgoLined
\KwIn{Data $(X, y)$, hyperparams $\{T, T_w, C, \rho_0, p\}$}
\KwOut{Feature ranking by $|\bg|$}
Init $\bg,\by,\bz \leftarrow \mathbf{1}$;\quad $\bv,\bw \leftarrow \mathbf{0}$;\quad $\rho \leftarrow \rho_0$\;
\For{$t = 0$ \KwTo $T-1$}{
    \If{$t < T_w$}{
        \tcp{Phase 1: Warm-up}
        Adam step on $\mathcal{L}(\theta, \bg)$ with feat.~dropout\;
    }
    \Else{
        \tcp{Phase 2: Three-variable ADMM}
        $s_j \leftarrow \lVert\bW_1[:,j]\rVert_2 + \epsilon$\;
        \textbf{g-step:} Adam on $\mathcal{L}_\rho$ via~\eqref{eq:g-practical}\;
        $c \leftarrow \lVert\bz\rVert_1$;\quad
        $\bd \leftarrow \bg + \bv/\rho$\;
        \textbf{y-step:} Solve cubic~\eqref{eq:cubic} for $\tau$;\quad
        $\by \leftarrow \tau\,\bd$\;
        \textbf{z-step:} $\bz \leftarrow S(\bg + \bw/\rho,\; C/(N\rho\lVert\by\rVert_2))$\;
        \textbf{dual:} $\bv \leftarrow \bv + \rho(\bg - \by)$;\quad
        $\bw \leftarrow \bw + \rho(\bg - \bz)$\;
        \If{$t \bmod 5 = 0$}{
            Adapt $\rho$ (residual balancing)\;
            $\bv,\bw \leftarrow \bv,\bw \cdot \rho_\text{old}/\rho_\text{new}$\;
        }
    }
}
\Return ranking by $|g_j|$ descending\;
\end{algorithm}

% ================================================================
%  4. EXPERIMENTS
% ================================================================
\section{Experiments}\label{sec:exp}

We evaluate our method on three benchmarks of increasing complexity: synthetic datasets with known informative features, a causal DAG dataset with two-level ground truth, and the MADELON real-world challenge dataset.

\subsection{Synthetic Benchmarks}\label{sec:synth}

\paragraph{Setup.}
We use four synthetic tasks from the Feature-Selection-Benchmark suite \citep{guyon2003introduction}:
\textbf{Ring} ($k{=}2$, radial band), \textbf{XOR} ($k{=}2$, quadrant interaction),
\textbf{Ring+XOR} ($k{=}4$), and \textbf{Ring+XOR+Sum} ($k{=}6$, with additive noise).
Samples are drawn uniformly in $[0,1]^m$ with $n{=}1{,}000$, and feature dimensions range over $m \in \{8, 16, 32, \ldots, 2048\}$.
The signal-to-noise ratio $k/m$ decreases from 25\% ($m{=}8$) to 0.1\% ($m{=}2048$).
All methods are evaluated via 6-fold CV with random feature permutation per fold.
The metric is \textbf{best-$k$}: the fraction of true informative features appearing in the top-$k$ ranked features (Precision@$k$).

\paragraph{Baselines.}
We compare against 11 methods spanning filter (MI, mRMR, Relief), tree-based (RF, TreeSHAP), and neural (LassoNet, CancelOut-sigmoid, CancelOut-softmax, CAE, FSNet, DeepPINK) categories, all run on the same data splits with the same evaluation protocol.

\paragraph{Results.}
Table~\ref{tab:synthetic} reports average best-$k$ across all values of $m$.
Our method achieves 66.8\% overall, ranking \textbf{first among all neural network methods} (vs.\ LassoNet 56.4\%, CancelOut-sig 47.0\%, DeepPINK 31.2\%) and \textbf{third overall} behind RF (82.2\%) and TreeSHAP (75.1\%).
Notably, on the \textbf{XOR} task---which requires detecting feature interactions undetectable by marginal statistics---our method achieves 83.3\%, substantially surpassing all non-neural methods (RF 54.5\%, mRMR 11.4\%, MI 18.2\%).
This confirms the advantage of neural gating for capturing nonlinear dependencies.

At $m{=}128$, where the signal-to-noise ratio is 1.6\% for XOR, our method achieves perfect 100\% best-$k$ accuracy, matching LassoNet (100.0\%) while TreeSHAP collapses to chance (49.8\%)---confirming that tree-based importance scores fail to detect interaction features at moderate dimensionality.
Performance degrades gracefully up to $m{=}256$ (100\% for XOR) and begins declining at $m{=}512$, indicating room for improvement in the ultra-high-dimensional regime.
Conversely, on Ring, TreeSHAP maintains near-perfect accuracy across all dimensions (99.2\% even at $m{=}2{,}048$), while both our method and LassoNet degrade substantially (Tables~\ref{tab:xor-dim}--\ref{tab:rxs-dim}).

\begin{table}[t]
    \centering
    \caption{Synthetic Datasets: Average Best-$k$ (\%, 6-fold CV) Across All $m$.}\label{tab:synthetic}
    \begin{tabular}{lcccc|c}
        \toprule
        \textbf{Method} & \textbf{Ring} & \textbf{XOR} & \textbf{R+X} & \textbf{R+X+S} & \textbf{Avg} \\
        \midrule
        RF          & \textbf{100.0} & 54.5 & \textbf{88.8} & \textbf{85.3} & \textbf{82.2} \\
        TreeSHAP    & \textbf{100.0} & 40.2 & 81.7 & 78.3 & 75.1 \\
        mRMR        & \textbf{100.0} & 11.4 & 81.7 & 74.7 & 67.0 \\
        \textbf{Ours}  & 63.9 & \textbf{83.3} & 54.2 & 65.8 & \underline{66.8} \\
        LassoNet    & 34.8 & 81.8 & 44.6 & 64.2 & 56.4 \\
        Relief      & 40.2 & 72.7 & 37.1 & 43.9 & 48.5 \\
        MI          & 92.6 & 18.2 & 39.6 & 40.8 & 47.8 \\
        CancelOut-$\sigma$ & 34.1 & 60.6 & 37.5 & 55.8 & 47.0 \\
        DeepPINK    & 21.2 & 34.1 & 21.2 & 48.3 & 31.2 \\
        CancelOut-sm & 26.5 & 31.8 & 24.6 & 40.3 & 30.8 \\
        CAE         & 19.7 & 22.7 & 19.2 & 25.8 & 21.9 \\
        FSNet       & 20.5 & 18.2 & 21.2 & 25.3 & 21.3 \\
        \bottomrule
    \end{tabular}
\end{table}

\subsection{Causal DAG Benchmark}\label{sec:dag}

\paragraph{Setup.}
We generate a sparse Bayesian network with $m{=}2{,}000$ features and $n{=}1{,}000$ observations.
The label is determined by a causal path through the DAG, defining two levels of ground truth:
$k \approx 15$ \emph{chain features} (direct causal ancestors) and $k_2 \approx 316$ \emph{chain+fork features} (including indirect influences through fork nodes).
We evaluate four metrics: \textbf{bestK} (Precision@$k$), \textbf{best2K} (Precision@$2k$), \textbf{bestK2} (Precision@$k_2$), and \textbf{best2K2} (Precision@$2k_2$).

\paragraph{Results.}
Table~\ref{tab:dag} shows the results.
Our method obtains modest scores on chain features (bestK = 7.8\%), where TreeSHAP (28.6\%) and RF (21.4\%) dominate.
However, it achieves the \textbf{highest scores on fork features}: bestK2 = 16.6\% and best2K2 = 32.4\%, substantially exceeding TreeSHAP (13.4\% / 17.3\%) and all other methods.
This indicates that Linearized ADMM with Ratio Norm is particularly effective at discovering \emph{indirect causal influences}---features that affect the target through intermediate variables in the DAG rather than through direct causal chains.

We hypothesize that the Ratio Norm's scale invariance helps detect fork features, which typically have weaker individual correlations with the target but collectively carry substantial information.
The standard $\ell_1$ penalty, being magnitude-sensitive, would disproportionately penalize these weaker signals.

\begin{table}[t]
    \centering
    \caption{DAG Dataset: Feature Recovery (\%, $m{=}2{,}000$).}\label{tab:dag}
    \begin{tabular}{lcccc}
        \toprule
        \textbf{Method} & \textbf{bestK} & \textbf{best2K} & \textbf{bestK2} & \textbf{best2K2} \\
        \midrule
        TreeSHAP & \textbf{28.6} & \textbf{42.9} & 13.4 & 17.3 \\
        RF       & 21.4 & 40.5 & 12.8 & 17.9 \\
        mRMR     & 16.7 & 19.0 & 6.6  & 11.5 \\
        CancelOut-$\sigma$ & 16.7 & 19.0 & 9.7 & 14.8 \\
        LassoNet & 14.3 & 14.3 & 6.8  & 11.1 \\
        DeepPINK & 14.3 & 14.3 & 8.4  & 12.3 \\
        MI       & 14.3 & 14.3 & 7.0  & 10.3 \\
        Relief   & 14.3 & 14.3 & 6.8  & 9.5 \\
        \textbf{Ours} & 7.8 & 8.9 & \textbf{16.6} & \textbf{32.4} \\
        CancelOut-sm & 9.5 & 23.8 & 7.0 & 12.1 \\
        CAE      & 2.4  & 11.9 & 4.9  & 8.0 \\
        FSNet    & 0.0  & 0.0  & 3.9  & 9.3 \\
        \bottomrule
    \end{tabular}
\end{table}

\subsection{MADELON Real-World Benchmark}\label{sec:madelon}

\paragraph{Setup.}
MADELON is a binary classification dataset from the NIPS 2003 Feature Selection Challenge \citep{guyon2004result}, with 500 features (only 20 informative, 96\% decoys), 2{,}000 training and 600 validation samples.
We select the top-$k$ features ($k{=}20$), train a Random Forest (500 trees) on the selected subset, and report AUROC and AUPRC.

\paragraph{Results.}
Table~\ref{tab:madelon} shows our method achieves AUROC = 0.964, matching RF (0.965) and TreeSHAP (0.964) within 0.1\%, and ranking \textbf{first among neural methods} (vs.\ LassoNet 0.950).
This confirms the method's effectiveness on real-world data where the underlying structure is unknown.

\begin{table}[t]
    \centering
    \caption{MADELON: Downstream Classification (top-20 features).}\label{tab:madelon}
    \begin{tabular}{lcc}
        \toprule
        \textbf{Method} & \textbf{AUROC} & \textbf{AUPRC} \\
        \midrule
        RF          & \textbf{0.965} & \textbf{0.965} \\
        TreeSHAP    & 0.964 & 0.965 \\
        \textbf{Ours} & 0.964 & 0.964 \\
        LassoNet    & 0.950 & 0.950 \\
        Relief      & 0.941 & 0.944 \\
        MI          & 0.833 & 0.831 \\
        FSNet       & 0.752 & 0.770 \\
        DeepPINK    & 0.742 & 0.768 \\
        CancelOut-sm & 0.740 & 0.743 \\
        CAE         & 0.711 & 0.713 \\
        CancelOut-$\sigma$ & 0.697 & 0.696 \\
        mRMR        & 0.640 & 0.646 \\
        \bottomrule
    \end{tabular}
\end{table}

\subsection{Analysis}\label{sec:analysis}

\paragraph{Strengths.}
Our method excels in three settings:
(1)~\emph{Interaction detection}: XOR-type patterns where marginal statistics fail (83.3\% vs.\ RF 54.5\%).
(2)~\emph{Fork-feature recovery}: Indirect causal influences in DAGs (32.4\% vs.\ TreeSHAP 17.3\%).
(3)~\emph{Real-world with high decoy ratio}: MADELON (96\% noise features) where it matches tree-based oracles.

\paragraph{Limitations.}
Performance drops in the ultra-high-dimensional regime ($m > 512$) on synthetic data, likely because the gate vector dimensionality grows while the signal remains fixed.
On the DAG benchmark, chain-feature detection (bestK = 7.8\%) is substantially below TreeSHAP (28.6\%), suggesting that the Ratio Norm's global scale-invariance may smooth over strong local signals.
Future work on the DAG task includes hyperparameter tuning beyond the current defaults, and potentially combining chain-optimized and fork-optimized runs.

\paragraph{Role of feature dropout.}
The 70\% feature dropout is essential for high-dimensional performance.
Without it, the network can memorize training data through noise features, producing nearly flat gate vectors.
The dropout forces the network to rely on a small subset of features per mini-batch, naturally encouraging sparsity in the learned gate.

\paragraph{Role of the Ratio Norm.}
The scale-invariant Ratio Norm treats features with different gate magnitudes equally, unlike the $\ell_1$ penalty which over-penalizes large-magnitude gates.
This is particularly important in the adaptive-penalty setting ($\lambda_j = C/s_j$), where features connected to different-scale first-layer weights would otherwise experience vastly different effective penalties under $\ell_1$.

% ================================================================
%  5. CONCLUSION
% ================================================================
\section{Conclusion}\label{sec:conclusion}

We presented a neural network feature selection framework based on Linearized ADMM with Ratio Norm regularization.
By keeping the learnable gate within Adam's gradient flow and applying the structured sparsity penalty through an epoch-level proximal step, our method bridges the gap between neural network flexibility and principled sparsity enforcement.
On diverse benchmarks, it consistently ranks as the best neural network method for feature selection, with particular strengths in detecting nonlinear interactions and indirect causal influences.
Future directions include extending to group-structured feature selection, adaptive warm-up scheduling, and integration with knockoff-based false discovery rate control.

% ================================================================
%  BACK MATTER
% ================================================================

\begin{contributions}
    Omitted for anonymous review.
\end{contributions}

\begin{acknowledgements}
    Omitted for anonymous review.
\end{acknowledgements}

\bibliography{references}

% ================================================================
%  SUPPLEMENTARY MATERIAL
% ================================================================
\newpage
\onecolumn
\title{Linearized ADMM with Ratio Norm for Neural Network Feature Selection\\(Supplementary Material)}
\maketitle

\appendix

\section{Hyperparameter Settings}\label{app:hparams}

Table~\ref{tab:hparams} lists all hyperparameters used in the experiments.
These values were selected via preliminary experiments on the XOR ($m{=}128$) dataset and used unchanged across all benchmarks.

\begin{table}[h]
    \centering
    \caption{Hyperparameter Settings.}\label{tab:hparams}
    \begin{tabular}{lcc}
        \toprule
        \textbf{Parameter} & \textbf{Value} & \textbf{Role} \\
        \midrule
        Learning rate           & 0.005   & Adam step size \\
        $C$ (sparsity coeff.)   & 0.05    & Ratio Norm penalty strength \\
        Total epochs $T$        & 500     & Training duration \\
        Warm-up epochs $T_w$    & 120     & Phase 1 epochs (no sparsity) \\
        Feature dropout $p$     & 0.7     & Fraction of features dropped \\
        $\rho_0$                & 200.0   & Initial ADMM penalty \\
        Batch size              & 64      & Mini-batch size \\
        Weight decay            & $10^{-3}$ & Adam $\ell_2$ regularization \\
        Hidden dimensions       & $[32, 32]$ & MLP hidden layer sizes \\
        Activation              & Mish    & Hidden layer activation \\
        $\epsilon$ (score floor) & $10^{-8}$ & Prevents division by zero in $\lambda_j$ \\
        \bottomrule
    \end{tabular}
\end{table}

\section{Detailed Results by Dimension}\label{app:bydim}

Tables~\ref{tab:xor-dim}--\ref{tab:rxs-dim} show best-$k$ accuracy for each synthetic task as $m$ increases, comparing our method against TreeSHAP and LassoNet at matched dimensions.

\begin{table}[h]
    \centering
    \caption{XOR ($k{=}2$): Best-$k$ (\%) by Dimension.}\label{tab:xor-dim}
    \resizebox{\columnwidth}{!}{%
    \begin{tabular}{lccccccccccc}
        \toprule
        $m$ & 2 & 4 & 8 & 16 & 32 & 64 & 128 & 256 & 512 & 1024 & 2048 \\
        \midrule
        Ours & 100 & 100 & 100 & 100 & 100 & 100 & 100 & 100 & 66.7 & 33.3 & 16.7 \\
        TreeSHAP & --- & --- & 99.9 & --- & 91.6 & --- & 49.8 & --- & 50.6 & --- & 50.6 \\
        LassoNet & --- & --- & 100.0 & --- & 100.0 & --- & 100.0 & --- & 100.0 & --- & 49.6 \\
        \bottomrule
    \end{tabular}%
    }
\end{table}

\begin{table}[h]
    \centering
    \caption{Ring ($k{=}2$): Best-$k$ (\%) by Dimension.}\label{tab:ring-dim}
    \resizebox{\columnwidth}{!}{%
    \begin{tabular}{lcccccccccc}
        \toprule
        $m$ & 8 & 16 & 32 & 64 & 128 & 256 & 512 & 1024 & 2048 \\
        \midrule
        Ours & 100 & 100 & 100 & 91.7 & 50 & 50 & 50 & 16.7 & 16.7 \\
        TreeSHAP & 99.2 & --- & 98.9 & --- & 99.2 & --- & 99.1 & --- & 98.8 \\
        LassoNet & 85.9 & --- & 76.9 & --- & 51.2 & --- & 50.4 & --- & 48.9 \\
        \bottomrule
    \end{tabular}%
    }
\end{table}

\begin{table}[h]
    \centering
    \caption{Ring+XOR ($k{=}4$): Best-$k$ (\%) by Dimension.}\label{tab:rx-dim}
    \resizebox{\columnwidth}{!}{%
    \begin{tabular}{lcccccccccc}
        \toprule
        $m$ & 4 & 8 & 16 & 32 & 64 & 128 & 256 & 512 & 1024 & 2048 \\
        \midrule
        Ours & 100 & 75 & 66.7 & 62.5 & 54.2 & 66.7 & 62.5 & 29.2 & 16.7 & 8.3 \\
        TreeSHAP & --- & 94.4 & --- & 95.7 & --- & 92.0 & --- & 75.3 & --- & 71.4 \\
        LassoNet & --- & 79.3 & --- & 80.4 & --- & 81.9 & --- & 75.7 & --- & 49.3 \\
        \bottomrule
    \end{tabular}%
    }
\end{table}

\begin{table}[h]
    \centering
    \caption{Ring+XOR+Sum ($k{=}6$): Best-$k$ (\%) by Dimension.}\label{tab:rxs-dim}
    \resizebox{\columnwidth}{!}{%
    \begin{tabular}{lcccccccccc}
        \toprule
        $m$ & 6 & 8 & 16 & 32 & 64 & 128 & 256 & 512 & 1024 & 2048 \\
        \midrule
        Ours & 100 & 88.9 & 72.2 & 69.4 & 66.7 & 63.9 & 69.4 & 55.6 & 38.9 & 33.3 \\
        TreeSHAP & --- & 85.1 & --- & 82.9 & --- & 82.1 & --- & 80.4 & --- & 69.0 \\
        LassoNet & --- & 83.8 & --- & 80.4 & --- & 82.1 & --- & 71.5 & --- & 65.7 \\
        \bottomrule
    \end{tabular}%
    }
\end{table}

\section{Full Derivation of the Three-Variable ADMM}\label{app:derivation}

We provide the complete mathematical derivation of all ADMM subproblems for the Ratio Norm objective.

\subsection{Problem Setup}

The regularized objective is:
\begin{equation}
    L(\bg) = \mathcal{L}(\theta,\bg) + \frac{C}{N}\frac{\lVert\bg\rVert_1}{\lVert\bg\rVert_2}.
\end{equation}
Introducing auxiliary variables $\by = \bg$ and $\bz = \bg$, the augmented Lagrangian is:
\begin{equation}
    L_\rho = \mathcal{L}(\theta,\bg) + \frac{C}{N}\frac{\lVert\bz\rVert_1}{\lVert\by\rVert_2}
    + \frac{\rho}{2}\left\lVert\bg - \by + \frac{\bv}{\rho}\right\rVert^2
    + \frac{\rho}{2}\left\lVert\bg - \bz + \frac{\bw}{\rho}\right\rVert^2.
\end{equation}

\subsection{$\bg$-Subproblem}

Fixing $\by$, $\bz$, $\bv$, $\bw$:
\begin{equation}
    \bg^{(k+1)} = \argmin_{\bg}\;\mathcal{L}(\theta,\bg)
    + \frac{\rho}{2}\left\lVert\bg - \by^{(k)} + \frac{\bv^{(k)}}{\rho}\right\rVert^2
    + \frac{\rho}{2}\left\lVert\bg - \bz^{(k)} + \frac{\bw^{(k)}}{\rho}\right\rVert^2.
\end{equation}
Taking the gradient and setting to zero:
\begin{equation}
    \nabla_{\bg}\mathcal{L} + \rho\bigl(\bg - \by^{(k)} + \bv^{(k)}/\rho\bigr)
    + \rho\bigl(\bg - \bz^{(k)} + \bw^{(k)}/\rho\bigr) = \mathbf{0}.
\end{equation}
Solving for $\bg$:
\begin{equation}
    \bg^{(k+1)} = \frac{1}{2}\left(\by^{(k)} + \bz^{(k)}
    - \frac{\bv^{(k)}}{\rho} - \frac{\bw^{(k)}}{\rho}
    - \frac{1}{\rho}\nabla_{\bg}\mathcal{L}\big|_{\bg^{(k+1)}}\right).
\end{equation}
The gradient $\nabla_{\bg}\mathcal{L}|_{\bg^{(k+1)}}$ is implicit.
We approximate it via the SADMM bridge (Section~\ref{sec:sadmm}): setting $\rho = 1/\alpha_k$ where $\alpha_k$ is Adam's per-parameter step size, and using the bias-corrected momentum $\hat{m}_k$:
\begin{equation}
    \bg^{(k+1)} = \frac{1}{2}\bigl(\by^{(k)} + \bz^{(k)}
    - \alpha_k\bv^{(k)} - \alpha_k\bw^{(k)} - \alpha_k\hat{m}_k\bigr).
\end{equation}

\subsection{$\by$-Subproblem (Ratio Norm Denominator)}

With $c^{(k)} = \lVert\bz^{(k)}\rVert_1$ and $\bd^{(k)} = \bg^{(k+1)} + \bv^{(k)}/\rho$, we solve:
\begin{equation}
    \by^{(k+1)} = \argmin_{\by}\;\frac{Cc^{(k)}}{N\lVert\by\rVert_2}
    + \frac{\rho}{2}\lVert\by - \bd^{(k)}\rVert^2.
\end{equation}
Taking the derivative:
\begin{equation}
    -\frac{Cc^{(k)}}{N\lVert\by\rVert_2^3}\by + \rho(\by - \bd^{(k)}) = \mathbf{0}.
\end{equation}
Rearranging:
\begin{equation}
    \left(-\frac{Cc^{(k)}}{N\lVert\by\rVert_2^3} + \rho\right)\by = \rho\bd^{(k)}.
\end{equation}

\textbf{Case 1}: $\bd^{(k)} \neq \mathbf{0}$.
Let $\by = \tau\bd^{(k)}$ for $\tau > 0$.
Substituting $\lVert\by\rVert_2 = \tau\lVert\bd^{(k)}\rVert_2$ and dividing by $\rho\bd^{(k)}$:
\begin{equation}
    \tau - 1 + \frac{Cc^{(k)}}{N\rho\tau^2\lVert\bd^{(k)}\rVert_2^3}\cdot\frac{1}{\tau} = 0
\end{equation}
Multiplying through by $\tau^2$:
\begin{equation}
    \tau^3 - \tau^2 - D^{(k)} = 0,\quad D^{(k)} = \frac{Cc^{(k)}}{N\rho\lVert\bd^{(k)}\rVert_2^3}.
\end{equation}
By Cardano's formula, the unique positive root is:
\begin{equation}
    \tau^{(k)} = \frac{1}{3} + \frac{1}{3}\left(\mathcal{C}^{(k)} + \frac{1}{\mathcal{C}^{(k)}}\right),
    \quad \mathcal{C}^{(k)} = \sqrt[3]{\frac{27D^{(k)}+2+\sqrt{(27D^{(k)}+2)^2-4}}{2}}.
\end{equation}

\textbf{Case 2}: $\bd^{(k)} = \mathbf{0}$.
Set $\by^{(k+1)} = \be^{(k)}$, a random vector with $\lVert\be^{(k)}\rVert_2 = (Cc^{(k)}/(N\rho))^{1/3}$.

\subsection{$\bz$-Subproblem (Ratio Norm Numerator)}

With $\by^{(k+1)}$ fixed:
\begin{equation}
    \bz^{(k+1)} = \argmin_{\bz}\;\frac{C}{N\lVert\by^{(k+1)}\rVert_2}\lVert\bz\rVert_1
    + \frac{\rho}{2}\left\lVert\bg^{(k+1)} - \bz + \frac{\bw^{(k)}}{\rho}\right\rVert^2.
\end{equation}
Taking the subgradient:
\begin{equation}
    \frac{C}{N\lVert\by^{(k+1)}\rVert_2}\sgn(\bz)
    - \rho\left(\bg^{(k+1)} - \bz + \frac{\bw^{(k)}}{\rho}\right) = \mathbf{0}.
\end{equation}
Solving yields the soft-thresholding operator:
\begin{equation}
    z_j^{(k+1)} = \begin{cases}
    a_j - \kappa & \text{if } a_j > \kappa, \\
    0 & \text{if } |a_j| \leq \kappa, \\
    a_j + \kappa & \text{if } a_j < -\kappa,
    \end{cases}
\end{equation}
where $a_j = g_j^{(k+1)} + w_j^{(k)}/\rho$ and $\kappa = C/(N\rho\lVert\by^{(k+1)}\rVert_2)$.

\subsection{$\ell_1$ Degeneration (Proof of Proposition~\ref{prop:l1})}

Replacing the Ratio Norm by $\ell_1$: $L = \mathcal{L}(\bg) + (C/N)\lVert\bg\rVert_1$.
The $\by$ variable is no longer needed ($\ell_2$ denominator absent).
With $\bg = \bz$ and a single dual variable $\bw$, the augmented Lagrangian is:
\begin{equation}
    L_\rho = \mathcal{L}(\bg) + \frac{C}{N}\lVert\bz\rVert_1
    + \frac{\rho}{2}\left\lVert\bg - \bz + \frac{\bw}{\rho}\right\rVert^2.
\end{equation}
The $\bg$-step gives $\bg^{(k+1)} = \bz^{(k)} - \alpha_k\bw^{(k)} - \alpha_k\nabla\mathcal{L}(\bz^{(k)})$.
The $\bz$-step applies soft-thresholding with threshold $C\alpha_k/N$:
\begin{equation}
    z_j^{(k+1)} = S\bigl(g_j^{(k+1)} + \alpha_k w_j^{(k)},\; C\alpha_k/N\bigr),
\end{equation}
which is the classical proximal gradient / iterative soft-thresholding update.\hfill$\square$

\section{DAG Dataset Generation}\label{app:dag}

The DAG dataset is generated as follows:
\begin{enumerate}
    \item Sample a random DAG with $m{=}2{,}000$ nodes and edge density 0.004.
    \item Assign edge weights uniformly from $[-1, -0.5] \cup [0.5, 1]$.
    \item Generate $n{=}1{,}000$ observations via ancestral sampling with $\sigma{=}0.2$ Gaussian noise.
    \item Identify the target node and trace its causal ancestors:
    \begin{itemize}
        \item $k$ chain features: nodes on directed paths to the target.
        \item $k_2$ chain+fork features: chain features plus nodes sharing a common child with chain features.
    \end{itemize}
    \item Generate model-X Gaussian knockoff features $\tilde{X}$ matching the empirical covariance.
\end{enumerate}

\end{document}
